\chapter{Overview and HITL Framework}
\label{ch:overview}

This chapter introduces the Human-in-the-Loop (HITL) Requirements
Traceability framework, its purpose, and the five-layer traceability
model that underpins the Trial Balance review process documentation.

% =============================================================================
\section{Purpose Statement}
\label{sec:purpose}

The HITL Requirements Traceability specification serves as the
\textbf{human validation checkpoint} for the automated requirements
extraction and traceability process. It is designed to:

\begin{enumerate}
  \item Provide full transparency into how business requirements were
        derived from source documents
  \item Enable human reviewers to verify the accuracy of automated
        extractions
  \item Establish a clear audit trail from source to implementation
  \item Support regulatory compliance by demonstrating requirement
        provenance
  \item Facilitate change impact analysis when source documents are
        updated
\end{enumerate}

% =============================================================================
\section{Five-Layer Traceability Model}
\label{sec:five-layer-model}

The traceability framework is organised into five distinct layers,
each representing a transformation stage from raw source material
to validated implementation artefacts.

\begin{figure}[h]
\centering
\begin{tikzpicture}[
    node distance=0.5cm,
    layer/.style={rectangle, draw=sapgrey, thick, minimum width=12cm,
                  minimum height=1cm, font=\small\bfseries},
    arrow/.style={->, thick, >=stealth, sapgrey}
]
    \node[layer, fill=sourcegreen!15] (L1) {Layer 1: Source Documents (6 files, 137 MB)};
    \node[layer, fill=sourcegreen!10, below=0.3cm of L1] (L2) {Layer 2: Extracted Machine-Readable (7 files, 522 KB)};
    \node[layer, fill=reqpurple!15, below=0.3cm of L2] (L3) {Layer 3: Business Requirements (27 requirements)};
    \node[layer, fill=specblue!15, below=0.3cm of L3] (L4) {Layer 4: Specifications (21 elements)};
    \node[layer, fill=schemaorg!15, below=0.3cm of L4] (L5) {Layer 5: Schemas \& Validation (6 schemas)};
    
    \draw[arrow] (L1.south) -- (L2.north);
    \draw[arrow] (L2.south) -- (L3.north);
    \draw[arrow] (L3.south) -- (L4.north);
    \draw[arrow] (L4.south) -- (L5.north);
\end{tikzpicture}
\caption{Five-Layer Traceability Model with artefact counts}
\label{fig:five-layer-model}
\end{figure}

\subsection{Layer 1: Source Documents}

Original business documents in their native formats:
\begin{itemize}[noitemsep]
  \item Business case documents (\texttt{.doc})
  \item Detailed Operating Instructions (\texttt{.docx})
  \item BPMN process maps (\texttt{.pdf})
  \item Working spreadsheets (\texttt{.xlsm})
\end{itemize}

\subsection{Layer 2: Extracted Machine-Readable}

Automated extraction outputs in machine-processable formats:
\begin{itemize}[noitemsep]
  \item Markdown with YAML frontmatter (text content)
  \item JSON (structured BPMN elements)
  \item YAML (schemas and specifications)
  \item CSV (sample data)
\end{itemize}

\subsection{Layer 3: Business Requirements}

Structured requirements derived from source documents:
\begin{itemize}[noitemsep]
  \item Executive Summary requirements (4)
  \item Scope requirements (6)
  \item Process Step requirements (8)
  \item AI Augmentation requirements (6)
  \item BSS Risk Assessment requirements (3)
\end{itemize}

\subsection{Layer 4: Specifications}

Formal specification elements implementing requirements:
\begin{itemize}[noitemsep]
  \item Process Steps (8)
  \item Decision Points (5)
  \item AI Capabilities (6)
  \item Controls (3)
\end{itemize}

\subsection{Layer 5: Schemas \& Validation}

JSON Schemas validating specification implementation:
\begin{itemize}[noitemsep]
  \item \schemaref{variance-record} --- variance analysis records
  \item \schemaref{decision-point} --- decision point states
  \item \schemaref{entity-params} --- per-entity parameters
  \item \schemaref{tb-extract} --- trial balance extraction
  \item \schemaref{pl-extract} --- P\&L extraction
  \item \schemaref{commentary-draft} --- LLM commentary drafts
\end{itemize}

% =============================================================================
\section{Traceability Marker Reference}
\label{sec:marker-reference}

Throughout this specification, colour-coded markers indicate
traceability links across layers:

\begin{table}[h]
\centering
\caption{Traceability Marker Reference}
\label{tab:marker-reference}
\begin{tabularx}{\textwidth}{@{}C{3cm}L{5cm}L{6cm}@{}}
\toprule
\textbf{Marker} & \textbf{Meaning} & \textbf{Example} \\
\midrule
\srcref{XXX} & Source document reference & \srcref{TB-BC-001} = Trial Balance Business Case \\
\reqref{XXX} & Business requirement ID & \reqref{TB-REQ-ES01} = Executive Summary req \\
\specref{XXX} & Specification element & \specref{TB-STEP-003} = Process step \\
\schemaref{XXX} & Schema definition & \schemaref{variance-record} = JSON Schema \\
\gapref{XXX} & Identified gap & \gapref{G-001} = Traceability gap \\
\bottomrule
\end{tabularx}
\end{table}

% =============================================================================
\section{HITL Review Checkpoints}
\label{sec:review-checkpoints}

Six human review checkpoints must be validated before the traceability
chain is considered complete:

\begin{table}[h]
\centering
\caption{HITL Review Checkpoints}
\label{tab:review-checkpoints}
\begin{tabularx}{\textwidth}{@{}C{1cm}L{4cm}L{7cm}C{2cm}@{}}
\toprule
\textbf{ID} & \textbf{Checkpoint} & \textbf{Review Criteria} & \textbf{Status} \\
\midrule
H1 & Source Completeness & All source documents listed and extracted & Pending \\
H2 & Extraction Accuracy & Extracted text matches source content & Pending \\
H3 & Requirement Derivation & Requirements accurately reflect source intent & Pending \\
H4 & Traceability Links & All links are valid and bidirectional & Pending \\
H5 & Schema Alignment & Schemas support all requirement fields & Pending \\
H6 & Gap Analysis & No orphan requirements or missing mappings & Pending \\
\bottomrule
\end{tabularx}
\end{table}

\begin{hitlbox}{H1 -- Source Completeness}
Human reviewer must verify that all 6 source documents are correctly
catalogued with accurate file sizes, types, and metadata. Any missing
documents must be identified and the extraction process re-run.
\end{hitlbox}

% =============================================================================
\section{Scope and Coverage}
\label{sec:scope-coverage}

\subsection{In Scope}

\begin{itemize}[noitemsep]
  \item Trial Balance review process across 234 Legal Entities
  \item AI-augmented variance analysis and commentary generation
  \item Month-end close controls (M-3 to M+5 timeline)
  \item FORTM pack generation and submission
  \item BSS Risk Assessment integration points
\end{itemize}

\subsection{Out of Scope}

\begin{itemize}[noitemsep]
  \item Country IFRS reporting (separate process)
  \item SAP/PeopleSoft integration detail
  \item Real-time intraday controls
  \item Audit committee reporting (downstream consumer)
\end{itemize}

\subsection{Coverage Statistics}

\begin{table}[h]
\centering
\caption{Traceability Coverage Statistics}
\label{tab:coverage-statistics}
\begin{tabular}{@{}lrr@{}}
\toprule
\textbf{Metric} & \textbf{Count} & \textbf{Coverage} \\
\midrule
Source Documents & 6 & 100\% catalogued \\
Extracted Files & 7 & 100\% traced to source \\
Business Requirements & 27 & 100\% traced to source \\
Specifications & 21 & 100\% traced to requirements \\
Schemas & 6 & 100\% traced to specifications \\
Identified Gaps & 2 & Documented with owners \\
\bottomrule
\end{tabular}
\end{table}